\documentclass[../../main.tex]{subfiles}
\graphicspath{{\subfix{../../image/}}}

\begin{document}

\section{习题}\label{section:ch06-习题}

在以下习题中，我们假设各 PDE 的系数光滑且满足一致椭圆性条件。另设 $U\subset\mathbb{R}^{n}$ 总为开、有界集，边界 $\partial U$ 光滑。

\begin{enumerate}
\item 考虑带位势函数 $c$ 的Laplace方程：
\begin{align}
-\Delta u+cu=0,
\end{align}
以及散度结构方程：
\begin{align}
-\operatorname{div}(aDv)=0,
\end{align}
其中函数 $a$ 为正。
\begin{enumerate}[(a)]
\item 证明若 $u$ 求解~(*) 且 $w>0$ 也求解~(*)，则 $v\triangleq u/w$ 在 $a\triangleq w^{2}$ 时求解~(**)。
\item 反之，证明若 $v$ 求解~(**)，则 $u\triangleq va^{1/2}$ 在某个位势 $c$ 时求解~(*)。
\end{enumerate}

\item 设
\begin{align*}
Lu=-\sum_{i,j=1}^{n}(a_{ij}u_{x_{i}})_{x_{j}}+cu.
\end{align*}
证明存在常数 $\mu>0$，使相应双线性型 $B[\cdot,\cdot]$ 满足 Lax--Milgram 定理的假设，只要
\begin{align*}
c(x)\geqslant-\mu\qquad (x\in U).
\end{align*}

\item 函数 $u\in H^{2}_{0}(U)$ 是双调和方程边值问题
\begin{align}
\begin{cases}
\Delta^{2}u=f & \text{于}\ U,\\
u=\dfrac{\partial u}{\partial\nu}=0 & \text{于}\ \partial U
\end{cases}
\end{align}
的弱解，只要对一切 $v\in H^{2}_{0}(U)$ 有
\begin{align*}
\int_{U}\Delta u\,\Delta v\,\mathrm{d}x=\int_{U}fv\,\mathrm{d}x.
\end{align*}
给定 $f\in L^{2}(U)$，证明~(*) 存在唯一弱解。

\item 设 $U$ 连通。函数 $u\in H^{1}(U)$ 是 Neumann 问题
\begin{align}
\begin{cases}
-\Delta u=f & \text{于}\ U,\\
\dfrac{\partial u}{\partial\nu}=0 & \text{于}\ \partial U
\end{cases}
\end{align}
的弱解，若对一切 $v\in H^{1}(U)$ 有
\begin{align*}
\int_{U}Du\cdot Dv\,\mathrm{d}x=\int_{U}fv\,\mathrm{d}x.
\end{align*}
设 $f\in L^{2}(U)$。证明~(*) 有弱解当且仅当
\begin{align*}
\int_{U}f\,\mathrm{d}x=0.
\end{align*}

\item 说明如何定义 $u\in H^{1}(U)$ 为带 Robin 边界条件的Poisson方程
\begin{align*}
\begin{cases}
-\Delta u=f & \text{于}\ U,\\
u+\dfrac{\partial u}{\partial\nu}=0 & \text{于}\ \partial U
\end{cases}
\end{align*}
的弱解。讨论给定 $f\in L^{2}(U)$ 时弱解的存在性与唯一性。

\item 设 $U$ 连通，$\partial U$ 由两个不相交闭集 $\Gamma_{1}$ 与 $\Gamma_{2}$ 组成。定义 $u$ 为带混合 Dirichlet--Neumann 边界条件的Poisson方程
\begin{align*}
\begin{cases}
-\Delta u=f & \text{于}\ U,\\
u=0 & \text{于}\ \Gamma_{1},\\
\dfrac{\partial u}{\partial\nu}=0 & \text{于}\ \Gamma_{2}
\end{cases}
\end{align*}
的弱解意味着什么。讨论弱解的存在性与唯一性。

\item 设 $u\in H^{1}(\mathbb{R}^{n})$ 有紧支集，且是半线性 PDE
\begin{align*}
-\Delta u+c(u)=f
\end{align*}
于 $\mathbb{R}^{n}$ 中的弱解，其中 $f\in L^{2}(\mathbb{R}^{n})$，$c:\mathbb{R}\to\mathbb{R}$ 光滑，$c(0)=0$、$c'\geqslant 0$。另设 $c(u)\in L^{2}(\mathbb{R}^{n})$。导出估计
\begin{align*}
\|D^{2}u\|_{L^{2}}\leqslant C\|f\|_{L^{2}}.
\end{align*}
（提示：模仿\cref{theorem:ch06-内部H2正则} 的证明，但不用截断函数 $\zeta$。）

\item 设 $u$ 是一致椭圆方程
\begin{align*}
Lu=-\sum_{i,j=1}^{n}a_{ij}(x)u_{x_{i}x_{j}}=0
\end{align*}
于 $U$ 中的光滑解。假设系数有有界导数。令 $v\triangleq|Du|^{2}+\lambda u^{2}$，证明若 $\lambda$ 充分大，则
\begin{align*}
Lv\leqslant 0\quad\text{于}\ U.
\end{align*}
由此推出
\begin{align*}
\|Du\|_{L^{\infty}(U)}\leqslant C\left(\|Du\|_{L^{\infty}(\partial U)}+\|u\|_{L^{\infty}(\partial U)}\right).
\end{align*}

\item 设 $u$ 是 $Lu=-\sum_{i,j=1}^{n}a_{ij}u_{x_{i}x_{j}}=f$ 于 $U$ 中的光滑解，$u=0$ 于 $\partial U$，其中 $f$ 有界。固定 $x_{0}\in\partial U$。$x_{0}$ 处的闸函数是 $C^{2}$ 函数 $w$，使
\begin{align*}
Lw\geqslant 1\ \text{于}\ U,\qquad w(x_{0})=0,\qquad w\geqslant 0\ \text{于}\ \partial U.
\end{align*}
证明若 $w$ 是 $x_{0}$ 处的闸函数，则存在常数 $C$ 使
\begin{align*}
|Du(x_{0})|\leqslant C\left|\frac{\partial w}{\partial\nu}(x_{0})\right|.
\end{align*}

\item 设 $U$ 连通。分别用 (a) 能量方法与 (b) 极值原理证明：Neumann 边值问题
\begin{align*}
\begin{cases}
-\Delta u=0 & \text{于}\ U,\\
\dfrac{\partial u}{\partial\nu}=0 & \text{于}\ \partial U
\end{cases}
\end{align*}
的唯一光滑解为 $u\equiv C$（某常数 $C$）。

\item 设 $u\in H^{1}(U)$ 是
\begin{align*}
-\sum_{i,j=1}^{n}(a_{ij}u_{x_{i}})_{x_{j}}=0
\end{align*}
于 $U$ 中的有界弱解。设 $\phi:\mathbb{R}\to\mathbb{R}$ 凸且光滑，令 $w=\phi(u)$。证明 $w$ 是弱下解；即对一切 $v\in H^{1}_{0}(U)$、$v\geqslant 0$，有 $B[w,v]\leqslant 0$。

\item 我们称一致椭圆算子
\begin{align*}
Lu=-\sum_{i,j=1}^{n}a_{ij}u_{x_{i}x_{j}}+\sum_{i=1}^{n}b_{i}u_{x_{i}}+cu
\end{align*}
满足弱极值原理，若对一切 $u\in C^{2}(U)\cap C(\overline{U})$，
\begin{align*}
\begin{cases}
Lu\leqslant 0 & \text{于}\ U,\\
u\leqslant 0 & \text{于}\ \partial U
\end{cases}
\end{align*}
蕴含 $u\leqslant 0$ 于 $U$。假设存在函数 $v\in C^{2}(U)\cap C(\overline{U})$，使 $Lv\geqslant 0$ 于 $U$ 且 $v>0$ 于 $\overline{U}$。证明 $L$ 满足弱极值原理。（提示：找一个没有零阶项的椭圆算子 $M$，使 $w\triangleq u/v$ 在区域 $\{u>0\}$ 中满足 $Mw\leqslant 0$。为此先计算 $(v^{2}w_{x_{i}})_{x_{j}}$。）

\item （Courant 极小极大原理）设 $L=-\sum_{i,j=1}^{n}(a_{ij}u_{x_{i}})_{x_{j}}$，其中 $((a_{ij}))$ 对称。假设算子 $L$（带零边界条件）有特征值 $0<\lambda_{1}<\lambda_{2}\leqslant\cdots$。证明
\begin{align*}
\lambda_{k}=\max_{S\in\Sigma_{k-1}}\min_{\substack{u\in S^{\perp}\\ \|u\|_{L^{2}}=1}}B[u,u]\qquad (k=1,2,\cdots).
\end{align*}
这里 $\Sigma_{k-1}$ 表示 $H^{1}_{0}(U)$ 的 $k-1$ 维子空间之集合。

\item 设 $\lambda_{1}$ 是一致椭圆、非对称算子
\begin{align*}
Lu=-\sum_{i,j=1}^{n}a_{ij}u_{x_{i}x_{j}}+\sum_{i=1}^{n}b_{i}u_{x_{i}}+cu
\end{align*}
（带零边界条件）的主特征值。证明“最大--最小”表示公式：
\begin{align*}
\lambda_{1}=\sup_{u}\inf_{x}\frac{Lu(x)}{u(x)},
\end{align*}
其中“sup”取遍 $u\in C^{\infty}(\overline{U})$、$u>0$ 于 $U$、$u=0$ 于 $\partial U$，“inf”取遍点 $x\in U$。（提示：考虑伴随算子 $L^{*}$ 对应于 $\lambda_{1}$ 的特征函数 $w_{1}^{*}$。）

\item （特征值与区域变化）考虑一族光滑、有界区域 $U(\tau)\subset\mathbb{R}^{n}$，光滑依赖于参数 $\tau\in\mathbb{R}$。当 $\tau$ 变化时，$\partial U(\tau)$ 上每点以速度 $v$ 运动。对每个 $\tau$，考虑特征值 $\lambda=\lambda(\tau)$ 及相应特征函数 $w=w(x,\tau)$：
\begin{align*}
\begin{cases}
-\Delta w=\lambda w & \text{于}\ U(\tau),\\
w=0 & \text{于}\ \partial U(\tau),
\end{cases}
\end{align*}
规范化使 $\|w\|_{L^{2}(U(\tau))}=1$。假设 $\lambda$ 与 $w$ 是 $\tau$ 与 $x$ 的光滑函数。证明 Hadamard 变分公式
\begin{align*}
\dot{\lambda}=-\int_{\partial U(\tau)}\left|\frac{\partial w}{\partial\nu}\right|^{2}v\cdot\nu\,\mathrm{d}S,
\end{align*}
其中 $\dot{\cdot}=\dfrac{\mathrm{d}}{\mathrm{d}\tau}$，$v\cdot\nu$ 为 $\partial U(\tau)$ 的法向速度。（提示：用\cref{section:appC-移动区域} 的微积分公式。）

\item （辐射条件）若我们分离变量以寻找波动方程的复值解，形式为 $u=e^{-\mathrm{i}\sigma t}w$（$w=w(x)$，$\sigma\in\mathbb{R}$、$\sigma\neq 0$），则导出特征值问题
\begin{align}
-\Delta w=\lambda w\quad\text{于}\ \mathbb{R}^{n},
\end{align}
其中 $\lambda\triangleq\sigma^{2}$。
\begin{enumerate}[(a)]
\item 证明 $w=e^{\mathrm{i}\sigma\omega\cdot x}$ 求解~(*)，只要 $|\omega|=1$。于是 $u=e^{\mathrm{i}\sigma(\omega\cdot x-t)}$ 是波动方程的行波解。
\item 证明对 $n=3$，函数 $\Phi\triangleq\dfrac{e^{\mathrm{i}\sigma|x|}}{4\pi|x|}$ 求解
\begin{align*}
-\Delta\Phi=\lambda\Phi+\delta_{0}\quad\text{于}\ \mathbb{R}^{3}.
\end{align*}
\item Sommerfeld 辐射条件要求~(*) 的解满足
\begin{align*}
\lim_{r\to\infty}r(w_{r}-i\sigma w)=0,
\end{align*}
其中 $w_{r}\triangleq Dw\cdot\dfrac{x}{|x|}$。证明 (a) 的解 $w$ 不满足此条件，但 (b) 的 $\Phi$ 满足。
\end{enumerate}

\item （续上题）证明若 $w$ 是 $\mathbb{R}^{3}$ 中特征值问题~(*) 的复值解且 $w$ 满足辐射条件，则 $w\equiv 0$。（提示：先注意
\begin{align*}
0=\int_{B(0,R)}\overline{w}\Delta w-w\Delta\overline{w}\,\mathrm{d}x=\int_{\partial B(0,R)}\overline{w}w_{r}-w\overline{w}_{r}\,\mathrm{d}S.
\end{align*}
利用此式与辐射条件证明
\begin{align*}
\int_{\partial B(0,R)}|w_{r}|^{2}+\sigma^{2}|w|^{2}\,\mathrm{d}S=\int_{\partial B(0,R)}|w_{r}-i\sigma w|^{2}\,\mathrm{d}S\to 0
\end{align*}
当 $R\to\infty$。现给定点 $x_{0}\in\mathbb{R}^{3}$，选 $R>|x_{0}|$。则
\begin{align*}
w(x_{0})=\int_{\partial B(0,R)}\Phi w_{r}-w\Phi_{r}\,\mathrm{d}S,
\end{align*}
其中 $\Phi=\Phi(x-x_{0})$。证明积分当 $R\to\infty$ 趋于零。）
\end{enumerate}

\end{document}
